\documentclass[10pt,a4paper]{article}
\usepackage{auditstyle}
\title{Fresh critical audit of the Glimm--Jaffe research program\\[3pt]
\large Severity-ranked proof check, source verification, and closure workflow}
\author{Audit target: \texttt{glimm\_jaffe/m6/main.tex}}
\date{9 August 2026}

\begin{document}
\maketitle

\begin{resultbox}
\textbf{Outcome first.}

\begin{itemize}
\item \textbf{Central Main Theorem 9.1: not refuted, but not certification-ready as
written.}  After the explicit local repairs recorded in this report and proved in the two
companion Lamport documents, its mathematical dependency chain closes conditionally on
the enumerated cutoff-construction and GJS inputs.
\item \textbf{Corollary 9.3 (``no symmetry breaking''): not proved as stated.}  Its only
cited internal theorem assumes a phase-decomposition hypothesis that Corollary 9.3 omits.
This is the audit's one confirmed proof-breaking issue.  It does not feed back into Main
Theorem 9.1.
\item \textbf{Citation correctness: three load-bearing pinpoint defects.}  Both FKN
citations to Simon Theorem III.12 are wrong; the hypercontractivity theorem numbers are
wrong; and the cited domain-smoothing theorem has a hypothesis not verified in the
manuscript.  Correct source theorems exist and repair all three imports.
\item \textbf{Local derivation defects: repairable.}  The manuscript has a false unit-cell
count, a false inequality for negative Euclidean time, an incomplete positivity choice in
the moment transfer, a dropped prefactor, an unrestricted spectral parameter, and omitted
signed truncation details.  None changes the qualitative theorem after correction.
\end{itemize}
\end{resultbox}

\tableofcontents

\section{Frozen audit object and standard of review}

The canonical target was fixed before the audit:
\[
 \begin{gathered}
 \texttt{glimm\_jaffe/m6/main.tex},\\
 \operatorname{SHA256}=\texttt{fdd67e696e68f93f5227ddfe32a6927a}\\[-2pt]
 \phantom{\operatorname{SHA256}=}\texttt{ace0099dc969819417ddcc9eca3f886a}.
 \end{gathered}
\]
The source has 1,022 lines.  A fresh \texttt{latexmk -pdf} build produced a 15-page PDF
with no unresolved references; layout warnings were limited to overfull boxes.  Successful
typesetting is not treated as mathematical verification.

The review standard was:

\begin{enumerate}[label=\textup{R\arabic*.}]
\item Every conclusion must be reachable from earlier displayed assumptions by an explicit
logical or analytic step.
\item Every use of an unbounded operator must state the relevant domain, closure, or
truncation passage.
\item Every limit must identify its topology and the theorem permitting the passage.
\item Every source citation must point to a result whose statement has the hypotheses and
conclusion actually used.
\item Contextual literature claims must not be promoted to mathematical inputs.
\item A defect is called \emph{proof-breaking} only if the claim does not follow from the
stated hypotheses.  A defect with a proved local correction is called \emph{repairable}.
\end{enumerate}

\section{Agentic workflow and gates}

The work was executed as a stateful audit workflow.  Each state had a testable exit gate;
failure returned the object to the preceding state rather than being concealed by prose.

\begin{longtable}{P{1.2cm}P{3.0cm}P{5.4cm}P{4.3cm}}
\toprule
State & Action & Concrete output & Exit gate\\
\midrule
\endhead
A0 & Freeze target & Canonical source, line count, SHA-256, fresh build & One unambiguous proof version\\
A1 & Reconstruct dependencies & Theorem-to-lemma graph and import ledger & Every main conclusion has predecessors\\
A2 & Local adversarial check & Recomputed constants, domains, signs, limits, spectral supports & Every internal arrow either passes or has a counterexample/repair\\
A3 & Source acquisition & All 15 bibliography entries present locally; three missing modern/method references downloaded afresh & Every citation has a local source artifact\\
A4 & Pinpoint source check & Exact theorem pages and scope comparison & Used statement is no stronger than source statement\\
A5 & Proof-break search & M3 hypothesis trace; historical and uniqueness scope audit & Unstated hypotheses isolated\\
A6 & Corrected reconstruction & Two independent Lamport-structured proofs & Every repair integrated and cross-referenced\\
A7 & Artifact QA & Compile, render every page, inspect every rendered page, clean temporary files & No TeX errors, clipping, blank pages, or unreadable evidence\\
\bottomrule
\end{longtable}

\section{Dependency graph and audit result}

The manuscript's central proof has the following actual dependency graph:
\[
\begin{array}{c}
\text{cutoff self-adjointness, simple positive ground state, FKN, domains}\\
\quad+\quad\text{GJS CE1, CE2, cutoff uniformity, scaling}
\end{array}
\Longrightarrow
\begin{array}{c}
(M_2)\ \text{uniform local moment}\\
(G)\ \text{uniform cutoff gap}
\end{array}
\]
\[
\begin{array}{c}
(M_2)+(G)+\text{finite propagation}+\text{patching}\\
+\text{affiliation}+\text{path regularity}
\end{array}
\Longrightarrow
\text{LPPL estimate}
\Longrightarrow
\begin{array}{c}
\text{full-net local-norm convergence}\cr
\text{local normality}
\end{array}
\]
\[
\text{local-norm convergence}+\text{patching}+\text{covariance}
\Longrightarrow
\begin{array}{c}
\text{spectral ground state}\cr
\text{translation and parity invariance.}
\end{array}
\]

The two GJS estimates are the only \emph{nonstandard weak-coupling estimates} in this
graph.  They are not the only logical inputs.  Therefore line 895's statement that all
conclusions are ``conditional only'' on GJS Theorems 1.1.7--1.1.8 is false as a literal
dependency claim.

\begin{longtable}{P{4.3cm}P{2.2cm}P{7.2cm}}
\toprule
Claim & Verdict & Reason\\
\midrule
Uniform local moment (M$_2$) & \Repair & Correct after choosing $F_1=F_2=|V(h)|$ and writing the unbounded truncation/weak-graph passage.\\
Uniform gap (G) & \Repair & CE2-to-semigroup argument and density argument pass after the cell-count, time-sign, moment-prefactor, and $\beta$ repairs.\\
Affine-path differentiability & \Repair & Argument passes with correct FKN/hypercontractivity citations and model-specific domain theorem.\\
Exact spectral filter & \Pass & Fourier sign, inverse-energy multiplier, integrability near $0$, and rapid tail all check.\\
Unbounded clustering & \Pass & Spectral supports, locality identity, affiliation, and Gaussian separator constants check.\\
One-path LPPL estimate & \Pass & Factor $32M$, distance loss, and shell count recompute correctly.\\
Full-net Cauchy limit/rate & \Pass & Cofinality and local dual completeness are used correctly.\\
Local normality & \Pass & Direct monotone-net argument proves norm-closedness in the required case.\\
Spectral ground-state passage & \Pass & The exact finite-volume spectral identity and the two-limit argument check.\\
Translation/$\mathbb Z_2$ invariance & \Pass & Covariance, cofinality, simplicity, and local norm passage check.\\
Canonical averaged vacuum & \Pass & The shifted plateaux contain $[-2n,2n]$ uniformly in the averaging parameter.\\
No symmetry breaking, Cor. 9.3 & \Fail & Companion M3 Theorem 4.2 assumes a broken-phase structure hypothesis absent from the corollary.\\
\bottomrule
\end{longtable}

\section{Confirmed proof-breaking issue}

\subsection{Corollary 9.3 omits the hypothesis used by M3}

The formal inference in the manuscript is invalid.

\LS{1}{1}{What Corollary 9.3 states.}

For an even polynomial in the weak-coupling region, it says the uniform gap is
incompatible with coexistence of two clustering phases exchanged by parity, citing
M3, Theorem 4.2.

\LS{1}{2}{What M3, Theorem 4.2 assumes.}

Lines 239--244 of \texttt{m3/m3.tex} impose Hypothesis 4.1: two translation-invariant
clustering ground states $\omega_\pm$, parity exchange, a nonzero local order parameter,
and---critically---identification of every parity-invariant ground state under discussion
with the symmetric midpoint $\tfrac12(\omega_++\omega_-)$.  The theorem begins
``Assume $P$ even and Hypothesis 4.1.''

\LS{1}{3}{Where the missing hypothesis is used.}

M3 first obtains a parity-invariant clustering limit $\omega$.  It then invokes the
hypothesis to write
\[
 \omega=\tfrac12(\omega_++\omega_-).
\]
For an order parameter with $\omega_\pm(A)=\pm m$, this gives
\[
 \omega(A)=0,
 \qquad
 \lim_{|a|\to\infty}\omega(A\tau_a(A))=m^2\ne0,
\]
contradicting clustering.  Without the midpoint identity, neither displayed equality
follows merely from parity invariance.

\LS{1}{4}{Logical diagnosis.}

M3 proves
\[
 \text{evenness}+\text{Hypothesis 4.1}+\text{uniform gap}
 \Longrightarrow\bot.
\]
Corollary 9.3 attempts to use this as
\[
 \text{evenness}+\text{uniform gap}+\text{two exchanged clustering states}
 \Longrightarrow\bot,
\]
which drops a premise.  This is the invalid inference
$(A\wedge B\Rightarrow C)\not\Rightarrow(A\Rightarrow C)$.

\LS{1}{5}{Exact repair choices.}

Either append all of M3 Hypothesis 4.1 to Corollary 9.3, or delete the corollary and
retain the proved parity invariance of the constructed cutoff-net limit.  The GJS
1975--1976 phase-transition papers do not, as cited in M6, automatically supply the exact
Hamiltonian phase-space classification needed here.

\section{Load-bearing source mismatches}

\subsection{FKN: Simon Theorem III.12 is the wrong theorem}

M6 lines 305 and 536 cite Simon Theorem III.12 for FKN.  The page image shows that
Theorem III.12 is an estimate for products of projections attached to Sobolev/Markov
regions.  It does not state either slab identity or the interacting semigroup formula.

The correct source chain is:
\[
 \begin{array}{ll}
 \text{Simon Theorem III.6} & \text{free-field FKN formula},\\
 \text{Simon Corollary V.14 and Theorem V.15} & \text{interacting FKN formula},\\
 \text{GRS Theorems II.14 and II.16} & \text{cutoff and general FKN formulas}.
 \end{array}
\]
The mathematical formula used in M6 is supported by these results, so this is a
repairable citation failure, not a refutation of the argument.

\subsection{Hypercontractivity: the exact theorem is Simon I.17/III.17}

M6 line 542 cites Simon Theorems V.10 and V.12 for
\[
 \|\mathsf P_TF\|_{p'}\leq\|F\|_p,
 \qquad\frac{p'-1}{p-1}=e^{2m_0T}.
\]
Simon Theorem I.17 is the free Nelson hypercontractivity theorem, and Theorem III.17
is its time-region form.  Theorem V.10 uses the earlier estimate inside a transfer-matrix
argument; Theorem V.12 is an essential-self-adjointness theorem.  The exponent arithmetic
in M6 is correct, but the pinpoint citation is not.

\subsection{Domain smoothing: use Expositions Theorem 7.4}

M6 line 214 cites Theorem 7.3 and Lemma 7.4 to infer
$\Ran e^{-tK(g)}\subset D(V(h))$.  Theorem 7.3 assumes the multiplication potential
$V$ itself is bounded below.  M6 does not establish that assumption for the cutoff Wick
interaction.  The next result, Expositions Theorem 7.4, is explicitly model-specific and
states
\[
 \Ran e^{-tH(g)}\subset D(H_0)\cap D(V')
\]
for every relevant polynomial multiplication $V'$.  Citing Theorem 7.4 repairs the
domain input exactly.

\section{Repairable internal defects with explicit derivations}

\subsection{False unit-cell bound}

M6 lines 376--377 use $L+1$ as an upper bound for the number of unit cells met by an
interval of length $L$.  Let $I=(-0.1,0.1)$; it has length $0.2$ and meets both
$[-1,0)$ and $[0,1)$ in positive measure, while $L+1=1.2<2$.  Thus the inequality is
false.

For half-open cells $[j,j+1)$, the correct integer bound is
\[
 N_I\leq\lfloor L\rfloor+2\leq L+2.                               \tag{6.1}
\]
Consequently the two displayed estimates must read
\begin{align*}
 \|Q\|_{\rm GJS}
 &\leq K_1K_2^{2r}r!(L+2)^{r/2}\prod_i\|f_i\|_2,\\
 \|Q^2\|_{\rm GJS}
 &\leq K_1K_2^{4r}(2r)!(L+2)^r\prod_i\|f_i\|_2^2.
\end{align*}
Only vector-dependent constants change; the decay exponent $m_1$ and gap proof do not.

\subsection{Wrong time-sign inequality}

M6 line 412 bounds $\|\psi_{T+s}\|$ by $\|\psi_T\|$ for both
$s=\pm t/2$.  Since $a\mapsto\|e^{-aH}\Om_0\|$ decreases, this holds for $s\geq0$
and reverses for $s<0$.

The corrected bound is
\[
 \frac{\|\psi_{T+s}\|}{\|\psi_T\|}
 \leq\begin{cases}1,&s\geq0,\\c_g^{-1},&s<0,
 \end{cases}
 \qquad c_g=\ip{\Om_g}{\Om_0}>0.                                  \tag{6.2}
\]
For fixed $g$ and fixed $s$, both ratios tend to $1$ as $T\to\infty$.  The finite
bound (6.2), not uniformity in $g$, is all that the weak-graph passage needs.

\subsection{Nonnegative FKN arguments at \texorpdfstring{$t=0$}{t=0}}

The slab lemma is stated for individually nonnegative $F_1,F_2$.  M6 line 352 specifies
only their product.  Set
\[
 F_1=F_2=|V(h)|.
\]
Then both hypotheses hold and $F_1F_2=V(h)^2$.  Bounded truncations
$|V(h)|\wedge M$ give the identity first; monotone convergence and CE1 give the
unbounded limit.

\subsection{Signed truncations in the gap transfer}

For a real polynomial multiplication operator $Q$, define
\[
 Q^{[M]}=(-M)\vee(Q\wedge M).
\]
Write $Q^{[M]}=(Q^{[M]})^+-(Q^{[M]})^-$ and apply the nonnegative FKN identity
bilinearly.  The two-time product is dominated by
\[
 |Q(-t/2)Q(t/2)|
 \leq\tfrac12\bigl(Q(-t/2)^2+Q(t/2)^2\bigr).                       \tag{6.3}
\]
CE1 bounds both terms.  On Hilbert space,
$Q^{[M]}\psi_a\to Q\psi_a$ follows from
$|Q-Q^{[M]}|^2|\psi_a|^2\leq Q^2|\psi_a|^2$.  These statements supply the
omitted limit theorem in lines 404--418.

\subsection{Dropped \texorpdfstring{$K_1$}{K1} and spectral range}

From the corrected moment bound
\[
 \E_{\mu_g}X^{2n}\leq C_8K_1(K')^{2n}(2n)!,
\]
one obtains
\[
 \E_{\mu_g}e^{s|X|}
 \leq2\max\{1,C_8K_1\}\sum_{n\geq0}(sK')^{2n},                    \tag{6.4}
\]
not the expression at line 455 with $K_1$ absent.  The radius $s<K'^{-1}$ is
unchanged.

In the final projection argument one must fix $0\leq\beta<m_1$.  For $\beta<0$,
$P_{[0,\beta]}\Om_g=0$, contrary to line 487.  Restricting $\beta$ gives
\[
 e^{-\beta t}\|P_{[0,\beta]}\psi_Q\|^2
 \leq\ip{\psi_Q}{e^{-tH(g)}\psi_Q}
 \leq C_Qe^{-m_1t},
\]
and the intended conclusion follows.

\subsection{Dimensionless weak-coupling reduction}

GJS page 590 supplies the dilation
\[
 (m_0,\lambda)\mapsto(s^{-1}m_0,s^{-2}\lambda).
\]
If the printed expansion is valid at an admissible $M_*$ for
$\lambda'<\lambda_*$, choose $s=m_0/M_*$.  Then
\[
 m_0'=M_*,\qquad \lambda'=\lambda M_*^2/m_0^2.
\]
Thus $\lambda/m_0^2<\lambda_*/M_*^2$ implies the printed hypotheses.  This explicit
reduction should accompany the definition
$\varepsilon_{\mathscr P}=\lambda_*/M_*^2$.

\section{Source-by-source cross-check}

All bibliography entries in M6 are present under \texttt{glimm\_jaffe/refs}; BDW,
RZT, and Lamport were downloaded afresh from their official author/arXiv locations.
The separate dossier contains the exact page images described below.

\begin{longtable}{P{2.0cm}P{4.2cm}P{5.5cm}P{2.5cm}}
\toprule
Source & Claim checked & Exact page/result inspected & Verdict\\
\midrule
\endhead
GJ I (1968) & cutoff self-adjointness; locality; cutoff independence in a light cone & article PDF pp. 5--6, especially locality/Trotter discussion & \Pass, broad citation\\
GJ II (1970) & ground state exists and is unique; locality/covariance & Thms. 2.2.1, 2.3.1; printed pp. 397--400 & \Pass\\
GJ III (1970) & averaged $w^*$ limit point and locally Fock/normal limit points & printed pp. 210--211, Thms. 2.1--2.3 & \Pass\\
GJ IV (1972) & completion of construction/Wightman functions & first-page Theorems A--B & Contextual pass\\
GJ (1971) & positivity/lower bounds/self-adjointness & Theorems 1--2, PDF pp. 2, 4 & Corroborating; not sufficient for every clause alone\\
Rosen (1970) & self-adjoint cutoff construction & Theorem 5.3 and following core argument, PDF pp. 22--23 & Corroborating\\
GJ Expositions & limit not proved & PDF pp. 54--56 (printed p. 50 and surrounding construction) & \Pass\\
GJ Expositions & finite propagation/patching & same chapter, PDF pp. 55--56 & \Pass, chapter citation broad\\
GJ Expositions & semigroup domain smoothing & Theorems 7.3--7.4, PDF pp. 165--167 & \Repair: cite Thm. 7.4\\
GJS (1974) & cutoff class and scaling & printed pp. 588--590 & \Pass\\
GJS (1974) & norm, degree class, CE2, CE1 & printed pp. 593--595, (1.1.5)--(1.1.10), Thms. 1.1.7--1.1.8 & \Pass\\
GJS (1974) & cutoff uniformity & printed p. 629 & \Pass for Thm. 1.1.8; CE2 already says independent of $h$\\
GRS (1975) & FKN & Thms. II.14 and II.16, Part I PDF pp. 39--40 & \Pass; use pinpoint citation\\
GRS (1975) & Problems 1--2 & printed pp. 125--126 and note added p. 128 & \Pass with historical qualification\\
Simon (1974) & free/interacting FKN & Thm. III.6; Cor. V.14; Thm. V.15 & \Repair\ citation\\
Simon (1974) & purported III.12 FKN citation & Thm. III.12, PDF p. 119 & \Fail as citation\\
Simon (1974) & hyper\-contractivity & Thm. I.17 and III.17 & \Repair\ citation\\
Simon (1974) & exponential stability; common core & Thms. V.7 and V.12 & \Pass in their actual roles\\
AHZ (1989) & Euclidean uniqueness/global Markov property in polynomial models & abstract and uniqueness conclusion near PDF p. 24 & Contextual pass; hypotheses matter\\
BDW (2025) & modern uniqueness & abstract: uniqueness of invariant measure for stochastic quantization above critical temperature & Scope mismatch if used for Hamiltonian uniqueness; M6 uses context only\\
RZT (2025) & Hamiltonian-renormalisation scope & abstract explicitly says finite volume & \Pass\\
GJS (1975--76) & phase transition/two pure phases & opening statements and theorem pages & Contextual pass; does not by itself supply M3 Hypothesis 4.1 as used\\
Lamport (2012) & hierarchical proof notation & worked examples and numbering discussion & Method followed in companion documents\\
\bottomrule
\end{longtable}

\section{Historical and scope corrections}

\subsection{GRS Problems 1--2}

GRS printed pages 125--126 formulate Euclidean local $L^p$ bounds and convergence of
finite-volume Gibbs states.  A note added on printed page 128 states that small-coupling
$p=1$ versions were solved by Newman.  M6 is justified in distinguishing its Hamiltonian
local-von-Neumann-algebra norm convergence from those Euclidean problems, but it should not
suggest without qualification that every formulation of Problems 1--2 remained open.

\subsection{Modern uniqueness citations}

BDW proves uniqueness of an invariant measure for the $\Phi^4_2$ stochastic quantization
equation above the critical temperature.  It does not itself prove uniqueness of locally
normal Hamiltonian ground states or convergence of the specific cutoff Hamiltonian net.
AHZ proves Euclidean Gibbs uniqueness/global Markov statements under its stated conditions.
M6 appropriately treats these as context, not inputs, but the wording should remain explicit.

\subsection{Exhaustive literature claim}

The sentence that Hamiltonian local-norm convergence ``has remained open in all regimes'' is
an exhaustive claim about the literature.  The finite list cited in M6 establishes that the
original construction used limit points and that RZT is finite volume; it does not establish
the absence of every later result in every regime.  Replace it by the defensible statement:
``The cited primary construction does not prove local-norm convergence of this cutoff net,
and no such theorem is identified in the references reviewed here.''

\section{Severity register}

\begin{longtable}{P{1.2cm}P{2.0cm}P{3.5cm}P{5.2cm}P{2.0cm}}
\toprule
ID & Severity & Location & Finding and mathematical consequence & Disposition\\
\midrule
\endhead
P1 & proof-breaking & Cor. 9.3; M3 lines 239--268 & Omitted phase-decomposition premise; corollary does not follow. & Add premise or delete\\
S1 & high & lines 305, 536 & Simon III.12 is not FKN. & Cite III.6, V.14--V.15 or GRS II.14/II.16\\
S2 & high & line 542 & V.10/V.12 are not the direct hypercontractivity theorem. & Cite I.17/III.17\\
S3 & high & line 214 & Theorem 7.3's bounded-below hypothesis is not checked. & Cite model-specific Thm. 7.4\\
L1 & moderate & lines 376--377 & $L+1$ cell count is false. & Replace by $L+2$\\
L2 & moderate & line 412 & Semigroup norm inequality fails for negative $s$. & Use sign split and $c_g^{-1}$\\
L3 & moderate & line 352 & Nonnegative inputs not individually specified. & Set both equal to $|V(h)|$\\
L4 & moderate & lines 404--418 & Signed/unbounded FKN limit under-described. & Use clipped truncations and (6.3)\\
L5 & low & line 455 & $K_1$ dropped from prefactor. & Restore it outside series\\
L6 & low & lines 484--489 & $\beta$ must be nonnegative. & State $0\leq\beta<m_1$\\
C1 & moderate & standing weak coupling & Dimensional reduction omitted. & Insert exact GJS scaling\\
C2 & moderate & line 895 & ``conditional only'' hides independent imports. & Enumerate all imports\\
H1 & moderate & introduction & ``open in all regimes'' not established by cited corpus. & Narrow claim\\
H2 & low & GRS discussion & Later p=1 note not mentioned. & Add qualification\\
\bottomrule
\end{longtable}

No P0 defect---a defect that survives all local repairs and breaks Main Theorem 9.1---was
found.  This negative finding is limited to the frozen source and imported theorems checked;
it is not a substitute for independent journal peer review.

\section{Closure plan for the research program}

The remaining work is finite and ordered by logical risk.

\begin{enumerate}[label=\textup{C\arabic*.}]
\item \textbf{Patch the theorem-bearing source.}  Apply S1--S3, L1--L6, and C1--C2
exactly as written in the two Lamport proofs.  Acceptance test: a line-by-line comparison
finds no use of the old constants, time-sign inequality, or theorem numbers.
\item \textbf{Resolve Corollary 9.3.}  Either state all of M3 Hypothesis 4.1 or remove the
corollary.  Acceptance test: every premise used in the midpoint contradiction appears in
the corollary statement.
\item \textbf{Narrow the literature claims.}  Incorporate the GRS note and replace the
exhaustive modern-open-problem assertion.  Acceptance test: every historical sentence is
entailed by a displayed primary-source page.
\item \textbf{Re-run the structured proof.}  Use the two companion Lamport documents as a
checklist against the patched M6 source.  Acceptance test: each $\langle1\rangle$ step maps
to a unique lemma/proposition in M6.
\item \textbf{Independent source audit of the only irreducible analytic import.}  Have a
second reader verify GJS (1.1.7)--(1.1.10), especially the localization-square convention
at sharp times and the cutoff-uniform scope.  Acceptance test: signed page references and
recomputed constants agree with the dossier.
\item \textbf{Certification build.}  Compile, run a reference/warning check, render every
page, and inspect every page.  Acceptance test: no unresolved references, clipped equations,
or unexplained external hypotheses.
\end{enumerate}

\section{Reference artifact manifest}

The full filenames are retained in \texttt{glimm\_jaffe/refs}.  The following SHA-256
values make the source audit reproducible; long titles are abbreviated only in this table.

\newcommand{\hashvalue}[2]{\texttt{#1}\newline\texttt{#2}}
\begin{longtable}{P{4.0cm}P{10.0cm}}
\toprule
Artifact & SHA-256\\
\midrule
\endhead
GJ I (1968) & \hashvalue{80993a0521b7f0de21b096c93365b3d8}{4cd569ed457dcb71b5766866b9c8ac07}\\
GJ II (1970) & \hashvalue{6c33fbc138766c344ee9eaca8bd78089}{3108a74549c8de74d995f08ec75f05d1}\\
GJ III (1970) & \hashvalue{8bfa327da9f3a78222243b3990b44693}{bd8e2cf1737c7a7d6f8749cfd3d9f961}\\
GJ IV (1972) & \hashvalue{b073f26cc15e08e66b4ea8480489109d}{7537f00e83ebfffd457cb62151483774}\\
GJ positivity (1971) & \hashvalue{655df0cadb54e132594d85a27b184971}{9128e8604fb4161f418a8d56d156ab97}\\
GJ Expositions & \hashvalue{d474345dc9336e968728e6ca2620d05f}{41ad21deb85143b48ccecb34b21410fe}\\
GJS (1974) & \hashvalue{e0a08c9438ff5c69f6d55f1c1b48e02}{01fc44f2313cef94ad4e87deaf872aa48}\\
GJS phase transitions (1975) & \hashvalue{dc20756fcdfa18992bb1034a35415d008}{e02339718beeca06d0ab74ec71c40e9}\\
GJS mean field I (1976) & \hashvalue{ef7957c5ef3eef3d867cf8db676bc9cb}{f03944d2f6ba79df3c68ca293fdc7b2e}\\
GJS mean field II (1976) & \hashvalue{7ffaff7052ea5f4dd1826fb1991f4eb6}{db9347275f150241ac6603cad2e47d64}\\
GRS Part I (1975) & \hashvalue{a97248302fe44eb9cddd690505c319f1}{bc9237e87f188fbefd74e4564e9945c5}\\
Rosen (1970) & \hashvalue{7336ed96d62107cd10459bd065199c06}{76c8b417fc9e74c450a20488f6c953ca}\\
Simon monograph & \hashvalue{06dc3239715cebb6a52ceec707e42bcc}{d0561ed2601370568d29790c3c59133c}\\
AHZ (1989) & \hashvalue{a4970d07869341427dd201a9fa397b8c}{8be825bb880234dfce61929882452165}\\
BDW arXiv:2504.08606 & \hashvalue{e1b9ebe4a8a702153e764b3b4a3e9b04}{e2adc9818725d0ebb3afeaf3cc954823}\\
RZT arXiv:2505.13030 & \hashvalue{50eaa5c031bcce2cf3ea057e2d6a9c3}{c793fe2734b4b88cae013d02fbf7cec23}\\
Lamport proof article & \hashvalue{c0fceb68ecfb1117138c39a942d61e3f}{17fb2f0c825067e18c9c657c393049a1}\\
\bottomrule
\end{longtable}

\section{Final audit verdict}

The highest-value correction is to withdraw or conditionalize Corollary 9.3.  The next
highest-value work is citation and domain repair, because the present pinpoint references
do not certify the formulas they are attached to.  Once those changes and the six local
derivation corrections are made, the central proof program has a complete conditional
route from the printed GJS estimates to the claimed cutoff-net limit.  The word
``conditional'' remains essential: the proof imports the cutoff construction, FKN,
semigroup smoothing, finite propagation, patching, covariance, and GJS itself.

\end{document}
